\section{Заключение}

В ходе работы был рассмотрен формат \texttt{bfloat16}, состоящий из бита
знака, восьми бит экспоненты и семи бит мантиссы. Реализованы функции
перевода десятичного числа в представление \texttt{bfloat16} и обратного
восстановления десятичного числа из строки из 16 бит.

При преобразовании в \texttt{bfloat16} число разбирается на целую и дробную
части, переводится в двоичный вид, нормализуется и округляется до семи бит
мантиссы. При обратном преобразовании из строки выделяются знак, экспонента
и мантисса, после чего по ним вычисляется исходное значение. Также
обрабатываются нуль, бесконечность и значение \texttt{NaN}.

